% !PRINTED 398
\weberpage{398}
From (7) it follows that when $u$ increases, $t$ must increase
as well, and hence that when $u$ has the smallest possible positive
value, $t$ is also as small as possible. One may therefore speak of a
smallest positive solution, which we shall denote by $T$, $U$. It is
obtained by taking $n$ in the formulae (4), (5) as small as possible,
that is, equal to the number of terms $\nu$ of the period. Hence if $\nu$
is even, only the equation (9), and not the equation (8), is
solvable; but if $\nu$ is odd, both (8) and (9) are solvable. Since
this can depend only on the discriminant $D$, and not on the particular
number $\omega$, it follows

\gesperrt{9. that for a given discriminant the various
periods contain either all an even or all
an odd number of terms.}

If $D$ is even, and hence divisible by 4, then $t$ must also be
even, and the equation (7) can be divided term by term by 4.
If $D$ is odd, then $t$, $u$ are likewise either both even
or both odd; if they are both odd, then their squares,
like all odd square numbers, are congruent to
1 modulo 8, and hence we must have $D \equiv 5\ (\mathrm{mod}\ 8)$. But it is not
always possible, when $D \equiv 5\ (\mathrm{mod}\ 8)$, to solve the equation (7)
in odd numbers\footnote{Cf. \gesperrt{Cayley}, Note sur l'équation $x^2 - D y^2 = \pm 4$, $D \equiv 5\ (\mathrm{mod}\ 8)$
Crelle's Journal, Bd. 53. Mathematical papers, Vol. IV. Tables for the
solutions of the Pell equation in the form $y^2 = a x^2 + 1$ have been
computed by \gesperrt{Degen} (Canon Pellianus, Havniae 1817). Such a table is also to be found in \gesperrt{Legendre}'s
number theory (German translation by \gesperrt{Maser}).}. If $D \equiv 1\ (\mathrm{mod}\ 8)$, then
$t$, $u$ must be even.

The examples of \S.\,126, treated in this way, yield
the following smallest positive solutions of the Pell
equation, where the factor 4 has been cancelled wherever possible.
\[
\begin{aligned}
5^2 - 29.1^2 &= -4\\
70^2 - 29.13^2 &= -1\\
170^2 - 19.39^2 &= 1\\
6^2 - 37.1^2 &= -1\\
35^2 - 34.6^2 &= 1.
\end{aligned}
\]
